%!TeX root=../gardentop.tex
\chapter{»I am Colin«} 
	
\begin{figure}[t!]
\centering
\includegraphics[width=\linewidth]{march}
\end{figure}

	\lettrine[lines=6]{M}{ary} took the picture back to the house when she went to her supper and she showed it to Martha.

\zz
»Eh!« said Martha with great pride. »I never knew our Dickon was as clever as that. That there's a picture of a missel thrush on her nest, as large as life an' twice as natural.«

\zz
Then Mary knew Dickon had meant the picture to be a message. He had meant that she might be sure he would keep her secret. Her garden was her nest and she was like a missel thrush. Oh, how she did like that queer, common boy!

She hoped he would come back the very next day and she fell asleep looking forward to the morning.

But you never know what the weather will do in Yorkshire, particularly in the springtime. She was awakened in the night by the sound of rain beating with heavy drops against her window. It was pouring down in torrents and the wind was »wuthering« round the corners and in the chimneys of the huge old house. Mary sat up in bed and felt miserable and angry.

»The rain is as contrary as I ever was,« she said. »It came because it knew I did not want it.«

She threw herself back on her pillow and buried her face. She did not cry, but she lay and hated the sound of the heavily beating rain, she hated the wind and its »wuthering.« She could not go to sleep again. The mournful sound kept her awake because she felt mournful herself. If she had felt happy it would probably have lulled her to sleep. How it »wuthered« and how the big rain-drops poured down and beat against the pane!

»It sounds just like a person lost on the moor and wandering on and on crying,« she said.

She had been lying awake turning from side to side for about an hour, when suddenly something made her sit up in bed and turn her head toward the door listening. She listened and she listened.

»It isn't the wind now,« she said in a loud whisper. »That isn't the wind. It is different. It is that crying I heard before.«

The door of her room was ajar and the sound came down the corridor, a far-off faint sound of fretful crying. She listened for a few minutes and each minute she became more and more sure. She felt as if she must find out what it was. It seemed even stranger than the secret garden and the buried key. Perhaps the fact that she was in a rebellious mood made her bold. She put her foot out of bed and stood on the floor.

»I am going to find out what it is,« she said. »Everybody is in bed and I don't care about Mrs~Medlock—I don't care!«

There was a candle by her bedside and she took it up and went softly out of the room. The corridor looked very long and dark, but she was too excited to mind that. She thought she remembered the corners she must turn to find the short corridor with the door covered with tapestry—the one Mrs~Medlock had come through the day she lost herself. The sound had come up that passage. So she went on with her dim light, almost feeling her way, her heart beating so loud that she fancied she could hear it. The far-off faint crying went on and led her. Sometimes it stopped for a moment or so and then began again. Was this the right corner to turn? She stopped and thought. Yes it was. Down this passage and then to the left, and then up two broad steps, and then to the right again. Yes, there was the tapestry door.

She pushed it open very gently and closed it behind her, and she stood in the corridor and could hear the crying quite plainly, though it was not loud. It was on the other side of the wall at her left and a few yards farther on there was a door. She could see a glimmer of light coming from beneath it. The Someone was crying in that room, and it was quite a young Someone.

So she walked to the door and pushed it open, and there she was standing in the room!

It was a big room with ancient, handsome furniture in it. There was a low fire glowing faintly on the hearth and a night light burning by the side of a carved four-posted bed hung with brocade, and on the bed was lying a boy, crying fretfully.

Mary wondered if she was in a real place or if she had fallen asleep again and was dreaming without knowing it.

The boy had a sharp, delicate face the colour of ivory and he seemed to have eyes too big for it. He had also a lot of hair which tumbled over his forehead in heavy locks and made his thin face seem smaller. He looked like a boy who had been ill, but he was crying more as if he were tired and cross than as if he were in pain.

Mary stood near the door with her candle in her hand, holding her breath. Then she crept across the room, and as she drew nearer the light attracted the boy's attention and he turned his head on his pillow and stared at her, his gray eyes opening so wide that they seemed immense.

»Who are you?« he said at last in a half-frightened whisper. »Are you a ghost?«

»No, I am not,« Mary answered, her own whisper sounding half frightened. »Are you one?«

He stared and stared and stared. Mary could not help noticing what strange eyes he had. They were agate gray and they looked too big for his face because they had black lashes all round them.

»No,« he replied after waiting a moment or so. »I am Colin.«

»Who is Colin?« she faltered.

»I am Colin Craven. Who are you?«

»I am Mary Lennox. Mr~Craven is my uncle.«

»He is my father,« said the boy.

»Your father!« gasped Mary. »No one ever told me he had a boy! Why didn't they?«

»Come here,« he said, still keeping his strange eyes fixed on her with an anxious expression.

She came close to the bed and he put out his hand and touched her.

»You are real, aren't you?« he said. »I have such real dreams very often. You might be one of them.«

Mary had slipped on a woolen wrapper before she left her room and she put a piece of it between his fingers.

»Rub that and see how thick and warm it is,« she said. »I will pinch you a little if you like, to show you how real I am. For a minute I thought you might be a dream too.«

»Where did you come from?« he asked.

»From my own room. The wind wuthered so I couldn't go to sleep and I heard some one crying and wanted to find out who it was. What were you crying for?«

»Because I couldn't go to sleep either and my head ached. Tell me your name again.«

»Mary Lennox. Did no one ever tell you I had come to live here?«

He was still fingering the fold of her wrapper, but he began to look a little more as if he believed in her reality.

»No,« he answered. »They daren't.«

»Why?« asked Mary.

»Because I should have been afraid you would see me. I won't let people see me and talk me over.«

»Why?« Mary asked again, feeling more mystified every moment.

»Because I am like this always, ill and having to lie down. My father won't let people talk me over either. The servants are not allowed to speak about me. If I live I may be a hunchback, but I shan't live. My father hates to think I may be like him.«

»Oh, what a queer house this is!« Mary said. »What a queer house! Everything is a kind of secret. Rooms are locked up and gardens are locked up—and you! Have you been locked up?«

»No. I stay in this room because I don't want to be moved out of it. It tires me too much.«

»Does your father come and see you?« Mary ventured.

»Sometimes. Generally when I am asleep. He doesn't want to see me.«

»Why?« Mary could not help asking again.

A sort of angry shadow passed over the boy's face.

»My mother died when I was born and it makes him wretched to look at me. He thinks I don't know, but I've heard people talking. He almost hates me.«

»He hates the garden, because she died,« said Mary half speaking to herself.

»What garden?« the boy asked.

»Oh! just—just a garden she used to like,« Mary stammered. »Have you been here always?«

»Nearly always. Sometimes I have been taken to places at the seaside, but I won't stay because people stare at me. I used to wear an iron thing to keep my back straight, but a grand doctor came from London to see me and said it was stupid. He told them to take it off and keep me out in the fresh air. I hate fresh air and I don't want to go out.«

»I didn't when first I came here,« said Mary. »Why do you keep looking at me like that?«

»Because of the dreams that are so real,« he answered rather fretfully. »Sometimes when I open my eyes I don't believe I'm awake.«

»We're both awake,« said Mary. She glanced round the room with its high ceiling and shadowy corners and dim firelight. »It looks quite like a dream, and it's the middle of the night, and everybody in the house is asleep—everybody but us. We are wide awake.«

»I don't want it to be a dream,« the boy said restlessly.

Mary thought of something all at once.

»If you don't like people to see you,« she began, »do you want me to go away?«

He still held the fold of her wrapper and he gave it a little pull.

»No,« he said. »I should be sure you were a dream if you went. If you are real, sit down on that big footstool and talk. I want to hear about you.«

Mary put down her candle on the table near the bed and sat down on the cushioned stool. She did not want to go away at all. She wanted to stay in the mysterious hidden-away room and talk to the mysterious boy.

»What do you want me to tell you?« she said.

He wanted to know how long she had been at Misselthwaite; he wanted to know which corridor her room was on; he wanted to know what she had been doing; if she disliked the moor as he disliked it; where she had lived before she came to Yorkshire. She answered all these questions and many more and he lay back on his pillow and listened. He made her tell him a great deal about India and about her voyage across the ocean. She found out that because he had been an invalid he had not learned things as other children had. One of his nurses had taught him to read when he was quite little and he was always reading and looking at pictures in splendid books.

Though his father rarely saw him when he was awake, he was given all sorts of wonderful things to amuse himself with. He never seemed to have been amused, however. He could have anything he asked for and was never made to do anything he did not like to do.

»Every one is obliged to do what pleases me,« he said indifferently. »It makes me ill to be angry. No one believes I shall live to grow up.«

He said it as if he was so accustomed to the idea that it had ceased to matter to him at all. He seemed to like the sound of Mary's voice. As she went on talking he listened in a drowsy, interested way. Once or twice she wondered if he were not gradually falling into a doze. But at last he asked a question which opened up a new subject.

»How old are you?« he asked.

»I am ten,« answered Mary, forgetting herself for the moment, »and so are you.«

»How do you know that?« he demanded in a surprised voice.

»Because when you were born the garden door was locked and the key was buried. And it has been locked for ten years.«

Colin half sat up, turning toward her, leaning on his elbows.

»What garden door was locked? Who did it? Where was the key buried?« he exclaimed as if he were suddenly very much interested.

»It—it was the garden Mr~Craven hates,« said Mary nervously. »He locked the door. No one—no one knew where he buried the key.«

»What sort of a garden is it?« Colin persisted eagerly.

»No one has been allowed to go into it for ten years,« was Mary's careful answer.

But it was too late to be careful. He was too much like herself. He too had had nothing to think about and the idea of a hidden garden attracted him as it had attracted her. He asked question after question. Where was it? Had she never looked for the door? Had she never asked the gardeners?

»They won't talk about it,« said Mary. »I think they have been told not to answer questions.«

»I would make them,« said Colin.

»Could you?« Mary faltered, beginning to feel frightened. If he could make people answer questions, who knew what might happen!

»Every one is obliged to please me. I told you that,« he said. »If I were to live, this place would sometime belong to me. They all know that. I would make them tell me.«

Mary had not known that she herself had been spoiled, but she could see quite plainly that this mysterious boy had been. He thought that the whole world belonged to him. How peculiar he was and how coolly he spoke of not living.

»Do you think you won't live?« she asked, partly because she was curious and partly in hope of making him forget the garden.

»I don't suppose I shall,« he answered as indifferently as he had spoken before. »Ever since I remember anything I have heard people say I shan't. At first they thought I was too little to understand and now they think I don't hear. But I do. My doctor is my father's cousin. He is quite poor and if I die he will have all Misselthwaite when my father is dead. I should think he wouldn't want me to live.«

»Do you want to live?« inquired Mary.

»No,« he answered, in a cross, tired fashion. »But I don't want to die. When I feel ill I lie here and think about it until I cry and cry.«

»I have heard you crying three times,« Mary said, »but I did not know who it was. Were you crying about that?« She did so want him to forget the garden.

»I dare say,« he answered. »Let us talk about something else. Talk about that garden. Don't you want to see it?«

»Yes,« answered Mary, in quite a low voice.

»I do,« he went on persistently. »I don't think I ever really wanted to see anything before, but I want to see that garden. I want the key dug up. I want the door unlocked. I would let them take me there in my chair. That would be getting fresh air. I am going to make them open the door.«

He had become quite excited and his strange eyes began to shine like stars and looked more immense than ever.

»They have to please me,« he said. »I will make them take me there and I will let you go, too.«

Mary's hands clutched each other. Everything would be spoiled—everything! Dickon would never come back. She would never again feel like a missel thrush with a safe-hidden nest.

»Oh, don't—don't—don't—don't do that!« she cried out.

He stared as if he thought she had gone crazy!

»Why?« he exclaimed. »You said you wanted to see it.«

»I do,« she answered almost with a sob in her throat, »but if you make them open the door and take you in like that it will never be a secret again.«

He leaned still farther forward.

»A secret,« he said. »What do you mean? Tell me.«

Mary's words almost tumbled over one another.

»You see—you see,« she panted, »if no one knows but ourselves—if there was a door, hidden somewhere under the ivy—if there was—and we could find it; and if we could slip through it together and shut it behind us, and no one knew any one was inside and we called it our garden and pretended that—that we were missel thrushes and it was our nest, and if we played there almost every day and dug and planted seeds and made it all come alive—«

»Is it dead?« he interrupted her.

»It soon will be if no one cares for it,« she went on. »The bulbs will live but the roses—«

He stopped her again as excited as she was herself.

»What are bulbs?« he put in quickly.

»They are daffodils and lilies and snowdrops. They are working in the earth now—pushing up pale green points because the spring is coming.«

»Is the spring coming?« he said. »What is it like? You don't see it in rooms if you are ill.«

»It is the sun shining on the rain and the rain falling on the sunshine, and things pushing up and working under the earth,« said Mary. »If the garden was a secret and we could get into it we could watch the things grow bigger every day, and see how many roses are alive. Don't you see? Oh, don't you see how much nicer it would be if it was a secret?«

He dropped back on his pillow and lay there with an odd expression on his face.

»I never had a secret,« he said, »except that one about not living to grow up. They don't know I know that, so it is a sort of secret. But I like this kind better.«

»If you won't make them take you to the garden,« pleaded Mary, »perhaps—I feel almost sure I can find out how to get in sometime. And then—if the doctor wants you to go out in your chair, and if you can always do what you want to do, perhaps—perhaps we might find some boy who would push you, and we could go alone and it would always be a secret garden.«

»I should—like—that,« he said very slowly, his eyes looking dreamy. »I should like that. I should not mind fresh air in a secret garden.«

Mary began to recover her breath and feel safer because the idea of keeping the secret seemed to please him. She felt almost sure that if she kept on talking and could make him see the garden in his mind as she had seen it he would like it so much that he could not bear to think that everybody might tramp into it when they chose.

»I'll tell you what I think it would be like, if we could go into it,« she said. »It has been shut up so long things have grown into a tangle perhaps.«

He lay quite still and listened while she went on talking about the roses which might have clambered from tree to tree and hung down—about the many birds which might have built their nests there because it was so safe. And then she told him about the robin and Ben Weatherstaff, and there was so much to tell about the robin and it was so easy and safe to talk about it that she ceased to feel afraid. The robin pleased him so much that he smiled until he looked almost beautiful, and at first Mary had thought that he was even plainer than herself, with his big eyes and heavy locks of hair.

»I did not know birds could be like that,« he said. »But if you stay in a room you never see things. What a lot of things you know. I feel as if you had been inside that garden.«

She did not know what to say, so she did not say anything. He evidently did not expect an answer and the next moment he gave her a surprise.

»I am going to let you look at something,« he said. »Do you see that rose-coloured silk curtain hanging on the wall over the mantel-piece?«

Mary had not noticed it before, but she looked up and saw it. It was a curtain of soft silk hanging over what seemed to be some picture.

»Yes,« she answered.

»There is a cord hanging from it,« said Colin. »Go and pull it.«

Mary got up, much mystified, and found the cord. When she pulled it the silk curtain ran back on rings and when it ran back it uncovered a picture. It was the picture of a girl with a laughing face. She had bright hair tied up with a blue ribbon and her gay, lovely eyes were exactly like Colin's unhappy ones, agate gray and looking twice as big as they really were because of the black lashes all round them.

»She is my mother,« said Colin complainingly. »I don't see why she died. Sometimes I hate her for doing it.«

»How queer!« said Mary.

»If she had lived I believe I should not have been ill always,« he grumbled. »I dare say I should have lived, too. And my father would not have hated to look at me. I dare say I should have had a strong back. Draw the curtain again.«

Mary did as she was told and returned to her footstool.

»She is much prettier than you,« she said, »but her eyes are just like yours—at least they are the same shape and colour. Why is the curtain drawn over her?«

He moved uncomfortably.

»I made them do it,« he said. »Sometimes I don't like to see her looking at me. She smiles too much when I am ill and miserable. Besides, she is mine and I don't want every one to see her.«

There were a few moments of silence and then Mary spoke.

»What would Mrs~Medlock do if she found out that I had been here?« she inquired.

»She would do as I told her to do,« he answered. »And I should tell her that I wanted you to come here and talk to me every day. I am glad you came.«

»So am I,« said Mary. »I will come as often as I can, but«—she hesitated—»I shall have to look every day for the garden door.«

»Yes, you must,« said Colin, »and you can tell me about it afterward.«

He lay thinking a few minutes, as he had done before, and then he spoke again.

»I think you shall be a secret, too,« he said. »I will not tell them until they find out. I can always send the nurse out of the room and say that I want to be by myself. Do you know Martha?«

»Yes, I know her very well,« said Mary. »She waits on me.«

He nodded his head toward the outer corridor.

»She is the one who is asleep in the other room. The nurse went away yesterday to stay all night with her sister and she always makes Martha attend to me when she wants to go out. Martha shall tell you when to come here.«

Then Mary understood Martha's troubled look when she had asked questions about the crying.

»Martha knew about you all the time?« she said.

»Yes; she often attends to me. The nurse likes to get away from me and then Martha comes.«

»I have been here a long time,« said Mary. »Shall I go away now? Your eyes look sleepy.«

»I wish I could go to sleep before you leave me,« he said rather shyly.

»Shut your eyes,« said Mary, drawing her footstool closer, »and I will do what my Ayah used to do in India. I will pat your hand and stroke it and sing something quite low.«

»I should like that perhaps,« he said drowsily.

Somehow she was sorry for him and did not want him to lie awake, so she leaned against the bed and began to stroke and pat his hand and sing a very low little chanting song in Hindustani.

»That is nice,« he said more drowsily still, and she went on chanting and stroking, but when she looked at him again his black lashes were lying close against his cheeks, for his eyes were shut and he was fast asleep. So she got up softly, took her candle and crept away without making a sound.